\chapter{Brain Aneurysm Coil Embolization}

\section*{What Is This Test or Procedure?}
Brain aneurysm coil embolization is an interventional neuroradiology procedure in which platinum coils are placed inside an aneurysm of a brain artery to block blood flow into it and prevent rupture.
\section*{Alternative Names}
Endovascular coiling; aneurysm coiling; coil embolization of a cerebral aneurysm
\section*{Principle}
Microcatheters are advanced through the arteries to the aneurysm, where detachable coils are packed to fill the sac. The coils promote blood clotting within the aneurysm and exclude it from the circulation, reducing the risk of hemorrhage.
\section*{History}
Endovascular coiling was pioneered by Guido Guglielmi in the 1990s with the development of detachable coils, and became a widely used alternative to surgical clipping after the ISAT trial in 2002.
\section*{Why This Test or Procedure Is Ordered}
Ordered for ruptured cerebral aneurysms that require urgent treatment to prevent rebleeding, and for unruptured aneurysms judged to carry a significant rupture risk, especially when surgery is high risk or not preferred.
\section*{Is This a Primary or Secondary Test or Procedure}
Primary therapeutic procedure for selected cerebral aneurysms.
\section*{How This Test or Procedure Is Performed}
Under fluoroscopic guidance and general anesthesia, a catheter is introduced through a groin or wrist artery and advanced into the brain's circulation. A microcatheter is navigated into the aneurysm and detachable coils are deployed to pack it, after which the catheters are removed.
\section*{What Preparation Is Needed}
Informed consent, imaging to define the aneurysm, and review of renal function and allergy history if contrast is used. Patients with unruptured aneurysms may receive antiplatelet therapy.
\section*{Contraindications and Precautions}
Contraindications include severe coagulopathy, contrast allergy that cannot be managed, and unfavorable aneurysm anatomy. Precautions include small neck aneurysms, which may be better treated by clipping, and renal impairment.
\section*{Risks}
Thromboembolism causing stroke, aneurysm perforation and hemorrhage, vessel injury, groin complications, and contrast nephropathy. Procedure-related morbidity is higher after subarachnoid hemorrhage.
\section*{What Happens After the Test or Procedure}
The patient is monitored closely, often in a neurointensive care setting after rupture, and receives antiplatelet therapy when indicated. Follow-up angiography checks for aneurysm recanalization.
\section*{Understanding Your Results}
Success is indicated by complete or near-complete occlusion of the aneurysm with no new neurologic deficit. Residual flow or coil compaction may require retreatment.
\section*{Normal Results}
A fully occluded aneurysm with no evidence of bleeding, thrombus, or vessel injury and no new neurologic symptoms.
\section*{Follow-up Tests and Procedures}
Serial imaging is performed to detect recurrence, and risk-factor control such as blood pressure and smoking cessation is emphasized to reduce future aneurysm risk.
\section*{References}
\begin{enumerate}
  \item MedlinePlus. Aneurysm in the brain - repair. U.S. National Library of Medicine; 2025.
  \item Current Medical Diagnosis and Treatment. McGraw-Hill; 2025.
\end{enumerate}

\clearpage
